% Generated draft: requires mathematical and source review.
\documentclass[11pt]{article}
\usepackage[margin=25mm]{geometry}
\usepackage{fontspec}
\setmainfont{Latin Modern Roman}
\usepackage{amsmath,amssymb,mathtools,unicode-math}
\setmathfont{Latin Modern Math}
\usepackage{enumitem,needspace,xurl,hyperref}
\setlength{\parindent}{0pt}
\setlength{\parskip}{5pt}
\setlength{\emergencystretch}{3em}
\newcommand{\N}{\mathbb{N}}
\newcommand{\Z}{\mathbb{Z}}
\newcommand{\Q}{\mathbb{Q}}
\newcommand{\R}{\mathbb{R}}
\newcommand{\C}{\mathbb{C}}
\newcommand{\F}{\mathbb{F}}
\newcommand{\A}{\mathbb{A}}
\newcommand{\PP}{\mathbb{P}}
\newcommand{\E}{\mathbb{E}}
\newcommand{\eps}{\varepsilon}
\newcommand{\dd}{\,\mathrm{d}}
\begin{document}
\section*{UnsolvedMath ID: 800626}
\textbf{Source problem code:} catalogue entry
\subsection{Definitions and mathematical statement}
For \(d\in\{1,2\}\), equip \(\R^d\) with its Euclidean topology, coordinatewise addition, and \(d\)-dimensional Lebesgue measure. A set is \emph{meager} if it is a countable union of nowhere dense sets. A set is \emph{null} if its Lebesgue outer measure is zero. For \(A,B\subseteq\R^d\), put
\[
A+B=\{a+b:a\in A,\ b\in B\}.
\]
A set \(A\subseteq\R^d\) is \emph{meager-additive} if \(A+M\) is meager for every meager set \(M\subseteq\R^d\). It is \emph{null-additive} if \(A+N\) is null for every null set \(N\subseteq\R^d\). These definitions impose no additional measurability, definability, or cardinality assumptions on \(A\).

The two questions supplied as Problem 55 are:
\begin{enumerate}
\item Is it true that, for every \(X,Y\subseteq\R\), if both \(X\) and \(Y\) are meager-additive in \(\R\), then \(X\times Y\) is meager-additive in \(\R^2\)? Explicitly, must
\[
\forall M\subseteq\R^2,\qquad
M\text{ meager}\ \Longrightarrow\ (X\times Y)+M\text{ meager}
\]
hold under these hypotheses?
\item Is it true that, for every \(X,Y\subseteq\R\), if both \(X\) and \(Y\) are null-additive in \(\R\), then \(X\times Y\) is null-additive in \(\R^2\)? Explicitly, must
\[
\forall N\subseteq\R^2,\qquad
N\text{ null}\ \Longrightarrow\ (X\times Y)+N\text{ null}
\]
hold under these hypotheses?
\end{enumerate}

The supplied background asserts that both corresponding product implications hold for subsets of \(\{0,1\}^{\mathbb N}\), with the product regarded as a subset of \(\{0,1\}^{\mathbb N}\times\{0,1\}^{\mathbb N}\). The conventional interpretation uses coordinatewise addition modulo \(2\), the product topology, and the fair-coin product probability measure, together with their product structures. The supplied excerpt does not explicitly specify these structures or define additivity; the definitions above and this interpretation therefore require confirmation against the original source. The registry labels the real-line questions open; this is a supplied status, not an independently verified current assessment.
\subsection{Short English statement}
Must the Cartesian product of any two meager-additive subsets of \(\R\) be meager-additive in \(\R^2\)? Ask the same question separately for two null-additive subsets. The supplied source reports affirmative analogues for Cantor space and presents the real-line cases as open.
\subsection{Sources}
\begin{itemize}
\item Local UnsolvedMath registry: permanent ID \(800626\), source code \texttt{catalogue entry}, title ``Selection Principles in Mathematics --- Problem 55 ([prods]).'' The supplied record belongs to AMR Open Problem Lists and has upstream revision \texttt{d371ee917ac9647fbb06cd6fc529deec6f3dda82}. Its supplied status is \texttt{open}.
\item Tsaban, \emph{Selection Principles and Special Sets of Reals: Open Problems} (2006), Problem 55. Supplied source URL: \url{https://arxiv.org/abs/math/0609601}. The registry attributes the problem to Tsaban and gives proposed year 2006; this does not establish original authorship of the question.
\item Supplied provenance reports access on 2026-10-16 and extraction by \texttt{arxiv\_latex}. This access date and the source assertions are supplied provenance, not independent verification performed for this draft.
\item Supplied attribution/context: \texttt{[prods]}. The corresponding bibliographic entry is absent from the supplied record.
\end{itemize}
\end{document}
